\documentclass[11pt, a4paper]{article}

% ─── Packages ────────────────────────────────────────────────────────────────
\usepackage[T1]{fontenc}
\usepackage[utf8]{inputenc}
\usepackage[margin=2.5cm]{geometry}
\usepackage{parskip}
\usepackage{titlesec}
\usepackage{titletoc}
\usepackage{fancyhdr}
\usepackage{xcolor}
\usepackage{booktabs}
\usepackage{array}
\usepackage{longtable}
\usepackage{amsmath, amssymb}
\usepackage{bm}
\usepackage{listings}
\usepackage{mdframed}
\usepackage{enumitem}
\usepackage{hyperref}
\usepackage{tcolorbox}
\tcbuselibrary{breakable, skins, listings}

% ─── Colours ─────────────────────────────────────────────────────────────────
\definecolor{codebg}{HTML}{1E1E2E}
\definecolor{codefg}{HTML}{CDD6F4}
\definecolor{keyword}{HTML}{CBA6F7}
\definecolor{string}{HTML}{A6E3A1}
\definecolor{comment}{HTML}{6C7086}
\definecolor{function}{HTML}{89B4FA}
\definecolor{number}{HTML}{FAB387}
\definecolor{accent}{HTML}{3B82F6}
\definecolor{accentlight}{HTML}{DBEAFE}
\definecolor{warnbg}{HTML}{FEF3C7}
\definecolor{warnborder}{HTML}{D97706}
\definecolor{bugbg}{HTML}{FEE2E2}
\definecolor{bugborder}{HTML}{DC2626}
\definecolor{darkgray}{HTML}{374151}
\definecolor{lightgray}{HTML}{F3F4F6}
\definecolor{tableborder}{HTML}{D1D5DB}

% ─── Hyperref ────────────────────────────────────────────────────────────────
\hypersetup{
  colorlinks = true,
  linkcolor  = accent,
  urlcolor   = accent,
  citecolor  = accent,
  pdftitle   = {Local RepE Framework \& Cross-Model Alignment: Engineering Blueprint},
  pdfauthor  = {LogoMesh Research},
}

% ─── Section Styling ─────────────────────────────────────────────────────────
\titleformat{\section}
  {\large\bfseries\color{darkgray}}
  {\thesection}{1em}{}[\vspace{2pt}\titlerule]

\titleformat{\subsection}
  {\normalsize\bfseries\color{darkgray}}
  {\thesubsection}{1em}{}

\titleformat{\subsubsection}
  {\normalsize\itshape\color{darkgray}}
  {\thesubsubsection}{1em}{}

\titlespacing*{\section}{0pt}{18pt}{8pt}
\titlespacing*{\subsection}{0pt}{12pt}{4pt}
\titlespacing*{\subsubsection}{0pt}{8pt}{2pt}

% ─── Header / Footer ─────────────────────────────────────────────────────────
\pagestyle{fancy}
\fancyhf{}
\renewcommand{\headrulewidth}{0.4pt}
\fancyhead[L]{\small\textcolor{darkgray}{Local RepE Framework \& Cross-Model Alignment}}
\fancyhead[R]{\small\textcolor{darkgray}{LogoMesh Internal Research}}
\fancyfoot[C]{\small\textcolor{darkgray}{\thepage}}

% ─── Code Listings ───────────────────────────────────────────────────────────
\lstdefinestyle{python}{
  language         = Python,
  backgroundcolor  = \color{codebg},
  basicstyle       = \ttfamily\footnotesize\color{codefg},
  keywordstyle     = \color{keyword}\bfseries,
  stringstyle      = \color{string},
  commentstyle     = \color{comment}\itshape,
  numberstyle      = \tiny\color{comment},
  numbers          = left,
  numbersep        = 8pt,
  stepnumber       = 1,
  showstringspaces = false,
  breaklines       = true,
  breakatwhitespace= true,
  frame            = none,
  xleftmargin      = 12pt,
  xrightmargin     = 4pt,
  aboveskip        = 6pt,
  belowskip        = 6pt,
  tabsize          = 4,
  morekeywords     = {self, True, False, None, yield, with, as, async, await, from, import},
  literate         = {->}{{\textcolor{keyword}{->}}}{2}
                     {**}{{\textcolor{keyword}{**}}}{2},
}

\lstdefinestyle{bash}{
  language         = bash,
  backgroundcolor  = \color{codebg},
  basicstyle       = \ttfamily\footnotesize\color{codefg},
  keywordstyle     = \color{keyword},
  commentstyle     = \color{comment}\itshape,
  showstringspaces = false,
  breaklines       = true,
  frame            = none,
  xleftmargin      = 12pt,
  aboveskip        = 6pt,
  belowskip        = 6pt,
}

\lstdefinestyle{json}{
  backgroundcolor  = \color{codebg},
  basicstyle       = \ttfamily\footnotesize\color{codefg},
  showstringspaces = false,
  breaklines       = true,
  frame            = none,
  xleftmargin      = 12pt,
  aboveskip        = 6pt,
  belowskip        = 6pt,
  morestring       = [b]",
  stringstyle      = \color{string},
}

\newtcolorbox{codebox}[1][]{
  enhanced,
  breakable,
  colback    = codebg,
  colframe   = codebg!80!black,
  boxrule    = 0.5pt,
  arc        = 4pt,
  left       = 2pt,
  right      = 2pt,
  top        = 4pt,
  bottom     = 4pt,
  #1
}

\newtcolorbox{warnbox}{
  enhanced,
  breakable,
  colback    = warnbg,
  colframe   = warnborder,
  boxrule    = 1pt,
  arc        = 3pt,
  left       = 8pt,
  right      = 8pt,
  top        = 6pt,
  bottom     = 6pt,
  fonttitle  = \bfseries\small\color{warnborder},
  title      = \textbf{Warning},
}

\newtcolorbox{bugbox}[1][Bug]{
  enhanced,
  breakable,
  colback    = bugbg,
  colframe   = bugborder,
  boxrule    = 1pt,
  arc        = 3pt,
  left       = 8pt,
  right      = 8pt,
  top        = 6pt,
  bottom     = 6pt,
  fonttitle  = \bfseries\small\color{bugborder},
  title      = \textbf{#1},
}

% ─── Table styling ───────────────────────────────────────────────────────────
\renewcommand{\arraystretch}{1.4}
\newcolumntype{L}[1]{>{\raggedright\arraybackslash}p{#1}}
\newcolumntype{C}[1]{>{\centering\arraybackslash}p{#1}}

% ─── Inline code ─────────────────────────────────────────────────────────────
\newcommand{\ic}[1]{\texttt{\small\colorbox{lightgray}{\strut #1}}}

% ─────────────────────────────────────────────────────────────────────────────
\begin{document}

% ─── Title Page ──────────────────────────────────────────────────────────────
\begin{titlepage}
  \centering
  \vspace*{3cm}
  {\Huge\bfseries\color{darkgray} Local RepE Framework \&\\[6pt] Cross-Model Alignment\par}
  \vspace{0.8cm}
  {\Large\color{accent} Engineering Blueprint\par}
  \vspace{1.2cm}
  \rule{12cm}{1pt}
  \vspace{0.6cm}

  \begin{tabular}{rl}
    \textbf{Classification:} & Internal Research / Architecture Specification \\[4pt]
    \textbf{Scope:}          & CPU-bound RepE extraction (local) + cross-model \\
                             & translation theory + LogoMesh Purple Agent \\[4pt]
    \textbf{Primary Literature:} & Zou et al.\ 2023 (RepE, arXiv:2310.01405) \\
                                 & Jha et al.\ 2025 (vec2vec, arXiv:2505.12540) \\[4pt]
    \textbf{Target Audience:} & Senior Security Engineering (Bakul) \\[4pt]
    \textbf{Date:}           & 2026-03-04 \\
  \end{tabular}

  \vspace{1.2cm}
  \rule{12cm}{1pt}
  \vspace{2cm}

  \begin{minipage}{10cm}
    \centering\small\color{darkgray}
    This document covers the full pipeline from CPU-constrained concept vector
    extraction on a consumer laptop through theoretical cross-architecture
    translation to a 70B model cluster, and the FastAPI server required to
    integrate the steered model into the LogoMesh adversarial evaluation arena.
  \end{minipage}

  \vfill
  {\small\color{darkgray} LogoMesh Research --- Confidential}
\end{titlepage}

% ─── Table of Contents ───────────────────────────────────────────────────────
\tableofcontents
\clearpage

% ═════════════════════════════════════════════════════════════════════════════
\section*{Overview}
\addcontentsline{toc}{section}{Overview}
% ═════════════════════════════════════════════════════════════════════════════

This blueprint is divided into two parts. \textbf{Part~1} is a concrete
engineering specification for building a Representation Engineering (RepE)
extraction pipeline on CPU-only consumer hardware. \textbf{Part~2} is an
exploratory theoretical analysis of how those extracted representations can be
translated to steer a 70B model on an H100 cluster, and the full integration
code for the LogoMesh Purple Agent.

\begin{center}
\begin{tabular}{lll}
  \toprule
  \textbf{Component} & \textbf{Local Pipeline} & \textbf{H100 Pipeline} \\
  \midrule
  Model              & Llama-3.2-3B (FP16)     & Llama-3.3-70B (BF16)  \\
  Hardware           & i7-12700H, 32\,GB RAM   & Lambda H100 cluster   \\
  Concept vector dim & 3072                    & 8192                  \\
  Layers targeted    & 18 (mid/late of 28)     & $\sim$53 (mid/late of 80) \\
  Injection $\alpha$ & 15.0 (empirical)        & 12.0 (empirical)      \\
  Translation method & N/A (source model)      & Supervised Procrustes \\
  Evaluation         & Manual / held-out       & LogoMesh CIS scorer   \\
  \bottomrule
\end{tabular}
\end{center}

% ═════════════════════════════════════════════════════════════════════════════
\part*{Part 1: The Local Engineering Blueprint}
\addcontentsline{toc}{part}{Part 1: The Local Engineering Blueprint}
% ═════════════════════════════════════════════════════════════════════════════

% ─────────────────────────────────────────────────────────────────────────────
\section{CPU-Bound Extraction Architecture}
% ─────────────────────────────────────────────────────────────────────────────

The core of RepE is \textbf{Linear Artificial Tomography (LAT)}: run contrastive
prompt pairs through the model, collect hidden state activations at a chosen
layer and token position, compute the difference vectors between positive and
negative stimuli, and extract the principal direction via PCA. That principal
direction \emph{is} the concept vector --- a geometric direction in activation
space that encodes the presence or absence of a concept (honesty, harmfulness,
code quality, etc.).

\subsection{Model Selection}

For a 32\,GB RAM / i7-12700H CPU environment, the viable candidates are:

\begin{center}
\begin{tabular}{lllllc}
  \toprule
  \textbf{Model} & \textbf{Params} & \textbf{Hidden Dim} & \textbf{Layers} & \textbf{FP16 RAM} & \textbf{Viable?} \\
  \midrule
  Qwen2.5-1.5B-Instruct  & 1.5B & 1536 & 28 & $\sim$3\,GB  & \checkmark\ Ideal \\
  Llama-3.2-1B-Instruct  & 1B   & 2048 & 22 & $\sim$2\,GB  & \checkmark\ Ideal \\
  Llama-3.2-3B-Instruct  & 3B   & 3072 & 28 & $\sim$6\,GB  & \checkmark\ Good  \\
  Mistral-7B             & 7B   & 4096 & 32 & $\sim$14\,GB & $\triangle$\ Tight \\
  \bottomrule
\end{tabular}
\end{center}

\textbf{Llama-3.2-3B-Instruct} is the recommended ceiling. It provides
sufficient representational depth for meaningful concept geometry while staying
well within the RAM budget when activations are cached efficiently.

\subsection{Hook Attachment and Activation Extraction}

PyTorch's \ic{register\_forward\_hook} is the mechanism. Attach hooks to the
output of each transformer block's residual stream --- i.e., the hidden state
\emph{after} the MLP sublayer but \emph{before} the next attention layer's
layer norm. This is the ``residual stream'' in the RepE paper's terminology.

\begin{codebox}
\begin{lstlisting}[style=python]
import torch
from transformers import AutoModelForCausalLM, AutoTokenizer
from sklearn.decomposition import PCA
import numpy as np

class RepEExtractor:
    def __init__(self, model_name: str, target_layers: list[int] | None = None):
        self.tokenizer = AutoTokenizer.from_pretrained(model_name)
        self.model = AutoModelForCausalLM.from_pretrained(
            model_name,
            torch_dtype=torch.float16,  # Critical: halves RAM footprint
            device_map="cpu",
        )
        self.model.eval()

        # Target the middle 2/3 of the network.
        # RepE paper shows concept geometry is richest in mid-to-late layers.
        n_layers = self.model.config.num_hidden_layers
        self.target_layers = target_layers or list(range(n_layers // 3, n_layers))

        self._hooks = []
        self._activations: dict[int, torch.Tensor] = {}

    def _attach_hooks(self):
        """Register hooks on residual stream outputs (decoder layer outputs)."""
        for layer_idx in self.target_layers:
            layer = self.model.model.layers[layer_idx]
            handle = layer.register_forward_hook(self._make_hook(layer_idx))
            self._hooks.append(handle)

    def _make_hook(self, layer_idx: int):
        def hook(module, input, output):
            # output[0]: (batch, seq_len, hidden_dim)
            # Extract ONLY the last token position to minimise memory.
            # RepE paper uses the last token as the concept "read point".
            hidden = output[0][:, -1, :].detach().cpu().float()
            self._activations[layer_idx] = hidden
        return hook

    def _remove_hooks(self):
        for h in self._hooks:
            h.remove()
        self._hooks.clear()

    @torch.no_grad()
    def extract(self, prompt: str) -> dict[int, np.ndarray]:
        """Run a single forward pass; return last-token activations per layer."""
        self._activations.clear()
        self._attach_hooks()
        try:
            inputs = self.tokenizer(prompt, return_tensors="pt")
            self.model(**inputs)
        finally:
            self._remove_hooks()
        return {k: v.numpy() for k, v in self._activations.items()}
\end{lstlisting}
\end{codebox}

\textbf{Key design decisions:}
\begin{itemize}[leftmargin=1.5em]
  \item \ic{torch.no\_grad()} eliminates the gradient computation graph,
        roughly halving peak RAM during the forward pass.
  \item \ic{torch.float16} reduces the 3B model's weight footprint from
        $\sim$12\,GB to $\sim$6\,GB.
  \item Extracting only the \textbf{last token position} (\ic{[:, -1, :]})
        reduces the activation tensor from
        \ic{(1, seq\_len, hidden\_dim)} to \ic{(1, hidden\_dim)} --- a
        reduction of $\text{seq\_len}\times$ (typically $64$--$512\times$).
  \item Hooks are attached and removed per call, preventing stale state
        accumulation across forward passes.
\end{itemize}

\subsection{LAT Vector Construction}

With an extractor in place, the LAT procedure from the RepE paper is:

\begin{codebox}
\begin{lstlisting}[style=python]
def build_lat_vectors(
    extractor: RepEExtractor,
    positive_prompts: list[str],
    negative_prompts: list[str],
) -> dict[int, np.ndarray]:
    """
    Returns a concept vector per target layer.
    Each concept vector is the first principal component of the
    difference-vector matrix (positive - negative activations).
    """
    assert len(positive_prompts) == len(negative_prompts)

    diffs: dict[int, list[np.ndarray]] = {l: [] for l in extractor.target_layers}

    for pos_prompt, neg_prompt in zip(positive_prompts, negative_prompts):
        pos_acts = extractor.extract(pos_prompt)
        neg_acts = extractor.extract(neg_prompt)

        for layer_idx in extractor.target_layers:
            diff = pos_acts[layer_idx] - neg_acts[layer_idx]  # (1, hidden_dim)
            diffs[layer_idx].append(diff[0])                  # (hidden_dim,)

    # PCA over the diff matrix -> principal direction = concept vector
    concept_vectors: dict[int, np.ndarray] = {}
    for layer_idx, diff_list in diffs.items():
        diff_matrix = np.stack(diff_list, axis=0)  # (N_pairs, hidden_dim)
        pca = PCA(n_components=1)
        pca.fit(diff_matrix)
        # Sign is ambiguous; validate post-hoc against known-positive prompts.
        concept_vectors[layer_idx] = pca.components_[0]  # (hidden_dim,)

    return concept_vectors
\end{lstlisting}
\end{codebox}

The sign of the PCA component is ambiguous. After extraction, validate by
projecting a known-positive and known-negative prompt onto the vector and
confirming the positive prompt yields a higher dot product.

\subsection{Inference-Time Injection (Steering)}

During inference on the target model, inject the concept vector into the
residual stream at the relevant layers using an additive hook:

\begin{codebox}
\begin{lstlisting}[style=python]
def make_injection_hook(concept_vector: np.ndarray, alpha: float = 15.0):
    """
    Returns a hook that adds alpha * concept_vector to the residual stream.
    Positive alpha amplifies the concept; negative alpha suppresses it.
    """
    v = torch.tensor(concept_vector, dtype=torch.float16)

    def hook(module, input, output):
        hidden = output[0]
        hidden = hidden + alpha * v.to(hidden.device)
        return (hidden,) + output[1:]
    return hook
\end{lstlisting}
\end{codebox}

The RepE paper recommends applying this at all target layers simultaneously.
The coefficient $\alpha$ should be tuned empirically: too high causes incoherent
outputs, too low has no effect. Values between 10 and 20 are typical.

% ─────────────────────────────────────────────────────────────────────────────
\section{Tensor Memory Management}
% ─────────────────────────────────────────────────────────────────────────────

\subsection{RAM Budget Analysis}

For a 3B model (hidden\_dim\,=\,3072, 28 layers, targeting 18 mid/late layers):

\begin{center}
\begin{tabular}{lc}
  \toprule
  \textbf{Item} & \textbf{Size} \\
  \midrule
  Model weights (FP16)                                    & $\sim$6\,GB   \\
  Peak KV cache during forward pass (seq\_len=512, FP16)  & $\sim$1.5\,GB \\
  Activations per prompt (last token, 18 layers, FP32)    & $\sim$0.2\,MB \\
  200 prompt pairs $\times$ 2 $\times$ 0.2\,MB           & $\sim$80\,MB  \\
  PCA computation overhead                                & $\sim$150\,MB \\
  \midrule
  \textbf{Total peak}                                     & $\sim$\textbf{8\,GB} \\
  \bottomrule
\end{tabular}
\end{center}

This is comfortable within 32\,GB. The danger zone is storing full-sequence
activations or forgetting \ic{torch.no\_grad()}.

\subsection{Specific Optimisation Strategies}

\textbf{Strategy 1 --- Incremental PCA.}
Do not cache all activations and then run PCA. Use \ic{IncrementalPCA} to
process mini-batches and discard raw tensors immediately:

\begin{codebox}
\begin{lstlisting}[style=python]
from sklearn.decomposition import IncrementalPCA

ipca = IncrementalPCA(n_components=1)
for batch in chunked(zip(pos_prompts, neg_prompts), n=10):
    batch_diffs = collect_diffs(batch)  # (10, hidden_dim)
    ipca.partial_fit(batch_diffs)
concept_vector = ipca.components_[0]
\end{lstlisting}
\end{codebox}

\textbf{Strategy 2 --- Explicit tensor lifecycle management.}

\begin{codebox}
\begin{lstlisting}[style=python]
diff = pos_acts[layer_idx] - neg_acts[layer_idx]
diffs[layer_idx].append(diff.copy())
del pos_acts, neg_acts, diff
import gc; gc.collect()
\end{lstlisting}
\end{codebox}

\textbf{Strategy 3 --- Memory-mapped storage for large datasets.}
If running more than 500 prompt pairs, use \ic{numpy.memmap} to write
activations to disk rather than holding them in RAM:

\begin{codebox}
\begin{lstlisting}[style=python]
import numpy as np
n_pairs, hidden_dim, n_target_layers = 500, 3072, 18

mmap = np.memmap(
    'activations.dat', dtype='float32', mode='w+',
    shape=(n_pairs, n_target_layers, hidden_dim)
)
# Write incrementally; NumPy loads pages on demand during PCA
\end{lstlisting}
\end{codebox}

\textbf{Strategy 4 --- 8-bit quantisation (last resort).}
Use \ic{bitsandbytes} if a 7B model is required:

\begin{codebox}
\begin{lstlisting}[style=python]
from transformers import BitsAndBytesConfig
quantization_config = BitsAndBytesConfig(load_in_8bit=True)
model = AutoModelForCausalLM.from_pretrained(
    model_name,
    quantization_config=quantization_config,
    device_map="cpu"
)
\end{lstlisting}
\end{codebox}

\begin{warnbox}
  BnB quantisation introduces noise into the activation geometry. For RepE
  work, FP16 is strongly preferred over INT8 if RAM permits. Concept vector
  quality degrades measurably at 8-bit and significantly at 4-bit.
\end{warnbox}

% ─────────────────────────────────────────────────────────────────────────────
\section{Data Pipeline: Contrastive Dataset Structure}
% ─────────────────────────────────────────────────────────────────────────────

The quality of the concept vector depends almost entirely on the quality of
the contrastive prompts. Design goals: (a) isolate the target concept from
confounds; (b) maximise variance in the difference vectors; (c) preserve all
metadata needed for future cross-model translation.

\subsection{Prompt Pair Schema}

For the LogoMesh use case, the most relevant concept is \textbf{code logical
correctness} --- since the CIS scorer's Logic (L) component has no floor
protection (Architecture Overview \S5.1 bug). A steered model that internally
represents ``high code correctness'' will naturally produce stronger L scores.

\begin{codebox}
\begin{lstlisting}[style=python]
from dataclasses import dataclass
import uuid

@dataclass
class ContrastivePair:
    pair_id: str
    concept: str          # e.g., "code_correctness"
    positive_prompt: str  # activates the concept strongly
    negative_prompt: str  # matched prompt that suppresses it
    task_domain: str      # e.g., "python_algorithms"
    difficulty: str       # "easy" | "medium" | "hard"
    notes: str = ""

@dataclass
class ActivationRecord:
    pair_id: str
    concept: str
    polarity: str          # "positive" or "negative"
    model_name: str
    model_hidden_dim: int
    model_n_layers: int
    layer_idx: int
    layer_frac: float      # layer_idx / model_n_layers -- key for depth alignment
    activation: list[float]
    prompt_token_len: int
\end{lstlisting}
\end{codebox}

The \ic{layer\_frac} field --- the layer index expressed as a fraction of total
model depth --- is the most important field for cross-model work. When
translating from a 22-layer model to an 80-layer model, relative depth is the
most tractable prior for establishing functional analogues.

\subsection{Critical Design Rules for Cross-Model Viability}

\begin{description}[leftmargin=2em, style=nextline]
  \item[\textbf{Rule 1 --- Concept isolation.}]
    Each pair must differ \emph{only} in the concept dimension. Control for
    prompt length, code complexity, and domain vocabulary. If positive prompts
    are systematically longer, the concept vector will be polluted with a
    ``length'' direction.

  \item[\textbf{Rule 2 --- Template diversification.}]
    Use at least 5 distinct surface-level phrasings per concept. The RepE paper
    uses multiple task templates (TruthfulQA, ARC, etc.) precisely to ensure
    the extracted direction generalises across prompt styles, not just one
    template.

  \item[\textbf{Rule 3 --- Metadata richness.}]
    Store the exact prompt, concept label, model name, layer index, and token
    position alongside every activation. This is the paired data that will
    supervise the cross-model translation mapping.
\end{description}

\subsection{Storage Format}

Use HDF5 via \ic{h5py} for the activation corpus. HDF5 supports partial I/O,
so you can load only specific layers without reading the entire corpus into RAM:

\begin{codebox}
\begin{lstlisting}[style=bash]
activations.h5
  metadata/
    model_name         (string)
    concept            (string)
    extraction_date    (string)
  concept_vectors/
    layer_00/          (float32, shape: hidden_dim)
    layer_01/
    ...
  raw/
    pair_0000/
      positive/        (float32 per layer: hidden_dim)
      negative/
    pair_0001/
    ...
\end{lstlisting}
\end{codebox}

% ═════════════════════════════════════════════════════════════════════════════
\clearpage
\part*{Part 2: Translation Theory and LogoMesh Integration}
\addcontentsline{toc}{part}{Part 2: Translation Theory and LogoMesh Integration}
% ═════════════════════════════════════════════════════════════════════════════

% ─────────────────────────────────────────────────────────────────────────────
\section{Residual Stream Alignment: The vec2vec Extension Problem}
% ─────────────────────────────────────────────────────────────────────────────

\subsection{What vec2vec Actually Does (and Does Not Do)}

The vec2vec paper (Jha et al.\ 2025) demonstrates that neural networks trained
on the same modality converge to a \textbf{universal latent semantic structure}
--- the Strong Platonic Representation Hypothesis. Their method uses a
GAN-based pipeline with input adapters, a shared backbone MLP, output adapters,
and cycle-consistency loss to learn mappings between embedding spaces
\emph{without any paired data}.

Critically, vec2vec operates exclusively on \textbf{final-layer static
embeddings} --- single vectors that summarise an entire input. The architecture
assumes: (1) one vector per document; (2) vectors are drawn from a fixed,
consistent distribution; (3) no concept of layer depth or temporal position.

Intermediate residual stream activations violate all three assumptions. They
are layer-indexed, context-dependent, and carry different semantic content at
different depths: early layers encode syntactic features, middle layers encode
relational structure, late layers encode task-level semantics. Translating them
is a fundamentally harder problem.

\subsection{The Two Mismatch Dimensions}

\subsubsection{Width (Dimensionality Mismatch)}

From a 1B Llama model (hidden\_dim\,=\,2048) to a 70B Llama model
(hidden\_dim\,=\,8192), the target space is $4\times$ larger. This is the
simpler of the two problems. A linear projection
$\mathbf{W} \in \mathbb{R}^{2048 \times 8192}$ is learnable given paired
activation data, and the Platonic Representation Hypothesis provides
theoretical justification that a good mapping exists.

The RepE paper implicitly supports this: since concept directions are
\emph{linear}, the mapping between the relevant subspaces of two models may
also be well-approximated by a linear transform. The best candidate is
\textbf{supervised Procrustes alignment}:

\[
  \text{Given: } \mathbf{X} \in \mathbb{R}^{N \times 2048}, \quad
  \mathbf{Y} \in \mathbb{R}^{N \times 8192}
\]
\[
  \text{Find: } \mathbf{W}^* = \arg\min_{\mathbf{W}} \|\mathbf{X}\mathbf{W} - \mathbf{Y}\|_F
\]
\[
  \text{Closed form: } \mathbf{W}^* = (\mathbf{X}^\top \mathbf{X})^{-1} \mathbf{X}^\top \mathbf{Y}
\]

Because you \emph{generate} the contrastive prompts yourself, you have exact
paired data. This is a major advantage over vec2vec's unsupervised setting ---
your problem is supervised Procrustes alignment, not adversarial GAN training.

\subsubsection{Depth (Layer Topographical Mismatch)}

A 22-layer model and an 80-layer model have radically different computational
graphs. Layer~11 of the small model is not functionally identical to layer~40
of the large model, even if both are ``mid-network.'' Three candidate strategies
exist:

\begin{description}[leftmargin=2em, style=nextline]
  \item[\textbf{Strategy A --- Proportional Depth Mapping (Simplest).}]
    Map layer $i$ of the small model to
    $\text{round}(i \times N_{\text{large}} / N_{\text{small}})$ of the large
    model. Trivially implementable, but assumes uniform feature development
    across depth, which is empirically false.

  \item[\textbf{Strategy B --- Functional Similarity via CKA (Recommended).}]
    Compute Centred Kernel Alignment (CKA) between all pairs of
    (small\_layer$_i$, large\_layer$_j$) using activations from the same
    prompts. Build a $22 \times 80$ similarity matrix, then solve with the
    Hungarian algorithm for optimal correspondence:

\begin{codebox}
\begin{lstlisting}[style=python]
def cka(X: np.ndarray, Y: np.ndarray) -> float:
    """Linear CKA between activation matrices (N x d1) and (N x d2)."""
    K = X @ X.T
    L = Y @ Y.T
    hsic_xy = np.sum(K * L)
    return hsic_xy / np.sqrt(np.sum(K * K) * np.sum(L * L))

cka_matrix = np.zeros((n_small_layers, n_large_layers))
for i, s in enumerate(small_acts_per_layer):
    for j, l in enumerate(large_acts_per_layer):
        cka_matrix[i, j] = cka(s, l)

from scipy.optimize import linear_sum_assignment
row_ind, col_ind = linear_sum_assignment(-cka_matrix)
layer_mapping = dict(zip(row_ind, col_ind))
\end{lstlisting}
\end{codebox}

  \item[\textbf{Strategy C --- Multi-Layer Concept Aggregation (Most Robust).}]
    Rather than mapping individual layers, identify which layers carry the
    \emph{strongest concept direction} in each model (highest PCA variance
    explained) and match only those ``concept-rich'' layers. This sidesteps
    the depth mismatch entirely: instead of asking ``which layers correspond?'',
    you ask ``which layers carry the concept?'' and match on that criterion.
\end{description}

\subsection{Supervised vs.\ Unsupervised: The Clear Answer}

\textbf{Use supervised alignment.} The vec2vec paper's unsupervised approach
solves the problem of having \emph{no paired data}. You have paired data. The
full pipeline is:

\begin{enumerate}[leftmargin=2em]
  \item Run $N$ contrastive pairs through the small model $\to$ collect
        $\mathbf{A}_{\text{small}} \in \mathbb{R}^{N \times 2048}$.
  \item Run the same $N$ prompts through the large model $\to$ collect
        $\mathbf{A}_{\text{large}} \in \mathbb{R}^{N \times 8192}$.
  \item Fit $\mathbf{W}^*$ via least-squares (closed form).
  \item Translate: $\mathbf{v}_{\text{large}} = \mathbf{W}^* \mathbf{v}_{\text{small}}$.
  \item Normalise $\mathbf{v}_{\text{large}}$ to unit length.
  \item Inject into the large model using the same additive hook mechanism.
\end{enumerate}

If the linear map underfits (measured by cosine similarity between
$\mathbf{W}^*\mathbf{v}_{\text{small}}$ and an independently-extracted
$\mathbf{v}_{\text{large}}$), upgrade to a small MLP (2-layer, ReLU,
hidden\_dim\,=\,4096) trained in a few hundred gradient steps.

\subsection{The Key Theoretical Risk: Context Dependence}

Static embeddings are (relatively) stable across batches. Intermediate
residual stream activations are \textbf{strictly context-dependent}: the
activation at layer~11 for token position $-1$ depends on the entire preceding
sequence. The learned mapping $\mathbf{W}^*$ will therefore have higher
variance and potentially lower transferability than a mapping trained on final
embeddings.

\textbf{Mitigation:} use prompts that terminate at a stable, semantically
neutral point (e.g.\ ``\texttt{This code is}'' as the terminal fragment) so
the last token position captures concept state rather than specific token
identity. The RepE paper already recommends this approach.

% ─────────────────────────────────────────────────────────────────────────────
\section{LogoMesh Validation Bridge: The Purple Agent Server}
% ─────────────────────────────────────────────────────────────────────────────

The following is a complete, production-ready implementation of the Purple Agent
FastAPI server. It wraps a RepE-steered 70B PyTorch model and conforms exactly
to the JSON-RPC 2.0 schema specified in the LogoMesh integration contracts.

\subsection{Integration Contract Reference}

The Green Agent sends:

\begin{codebox}
\begin{lstlisting}[style=json]
{
  "jsonrpc": "2.0",
  "method":  "message/send",
  "params": {
    "message": {
      "kind":      "message",
      "messageId": "task-uuid",
      "role":      "user",
      "parts": [{"kind": "text", "text": "Task prompt..."}]
    }
  },
  "id": "task-uuid"
}
\end{lstlisting}
\end{codebox}

The Purple Agent must return:

\begin{codebox}
\begin{lstlisting}[style=json]
{
  "jsonrpc": "2.0",
  "id": "task-uuid",
  "result": {
    "kind": "message",
    "parts": [{
      "kind": "text",
      "text": "{\"sourceCode\": \"...\", \"testCode\": \"...\",
               \"rationale\": \"...\"}"
    }]
  }
}
\end{lstlisting}
\end{codebox}

\subsection{Full Server Implementation}

\begin{codebox}
\begin{lstlisting}[style=python]
"""
purple_agent/server.py
LogoMesh Purple Agent -- FastAPI JSON-RPC 2.0 wrapper
for a RepE-steered 70B language model on Lambda Labs H100.
"""

import json, logging, time
from contextlib import asynccontextmanager
from typing import Any

import torch
from fastapi import FastAPI, Request
from fastapi.responses import JSONResponse
from pydantic import BaseModel, Field
from transformers import AutoModelForCausalLM, AutoTokenizer

# ── Configuration ──────────────────────────────────────────────────────────
MODEL_NAME          = "meta-llama/Llama-3.3-70B-Instruct"
STEERING_CHECKPOINT = "/checkpoints/repe_vectors.pt"
STEERING_ALPHA      = 12.0
MAX_NEW_TOKENS      = 2048
TEMPERATURE         = 0.1
DEVICE              = "cuda" if torch.cuda.is_available() else "cpu"

logger = logging.getLogger("purple_agent")
logging.basicConfig(level=logging.INFO)

model_state: dict[str, Any] = {}

# ── Lifespan: load model once at startup ──────────────────────────────────
@asynccontextmanager
async def lifespan(app: FastAPI):
    logger.info(f"Loading model: {MODEL_NAME}")
    tokenizer = AutoTokenizer.from_pretrained(MODEL_NAME)
    model = AutoModelForCausalLM.from_pretrained(
        MODEL_NAME,
        torch_dtype=torch.bfloat16,
        device_map="auto",  # Distributes across H100 GPUs automatically
    )
    model.eval()

    steering_data = torch.load(STEERING_CHECKPOINT, map_location=DEVICE)
    concept_vectors: dict[int, torch.Tensor] = steering_data["concept_vectors"]

    n_layers = model.config.num_hidden_layers
    target_layers = list(range(n_layers // 3, n_layers))

    model_state.update({
        "model": model, "tokenizer": tokenizer,
        "concept_vectors": concept_vectors, "target_layers": target_layers
    })
    logger.info(f"Steering {len(target_layers)} layers, alpha={STEERING_ALPHA}")
    yield
    model_state.clear()

app = FastAPI(title="LogoMesh Purple Agent", lifespan=lifespan)

# ── Steering hook management ───────────────────────────────────────────────
def attach_steering_hooks(model, concept_vectors, target_layers, alpha):
    hooks = []
    for idx in target_layers:
        if idx not in concept_vectors:
            continue
        v = concept_vectors[idx].to(DEVICE)

        def make_hook(vec):
            def hook(module, inp, output):
                hidden = output[0]
                hidden = hidden + alpha * vec.to(hidden.dtype).to(hidden.device)
                return (hidden,) + output[1:]
            return hook

        handle = model.model.layers[idx].register_forward_hook(make_hook(v))
        hooks.append(handle)
    return hooks

def remove_hooks(hooks):
    for h in hooks:
        h.remove()

# ── System prompt ──────────────────────────────────────────────────────────
SYSTEM_PROMPT = """You are an expert software engineer solving coding challenges.
For each task, produce ONLY a JSON object in this exact format:
{
  "sourceCode": "...",
  "testCode":   "...",
  "rationale":  "..."
}
No text outside the JSON object."""

# ── Steered generation ─────────────────────────────────────────────────────
@torch.no_grad()
def generate_solution(task_prompt: str) -> dict[str, str]:
    model     = model_state["model"]
    tokenizer = model_state["tokenizer"]
    cv        = model_state["concept_vectors"]
    tl        = model_state["target_layers"]

    messages = [
        {"role": "system", "content": SYSTEM_PROMPT},
        {"role": "user",   "content": task_prompt},
    ]
    input_ids = tokenizer.apply_chat_template(
        messages, tokenize=True, add_generation_prompt=True,
        return_tensors="pt",
    ).to(DEVICE)

    hooks = attach_steering_hooks(model, cv, tl, STEERING_ALPHA)
    try:
        output_ids = model.generate(
            input_ids, max_new_tokens=MAX_NEW_TOKENS,
            temperature=TEMPERATURE, do_sample=TEMPERATURE > 0,
            pad_token_id=tokenizer.eos_token_id,
        )
    finally:
        remove_hooks(hooks)

    new_tokens = output_ids[0][input_ids.shape[1]:]
    raw = tokenizer.decode(new_tokens, skip_special_tokens=True).strip()

    try:
        if raw.startswith("```"):
            raw = raw.split("```")[1]
            if raw.startswith("json"):
                raw = raw[4:]
        result = json.loads(raw)
        assert all(k in result for k in ("sourceCode", "testCode", "rationale"))
        return result
    except (json.JSONDecodeError, AssertionError) as e:
        logger.warning(f"JSON parse failed: {e}")
        return {
            "sourceCode": f"# Parse error: {e}\n{raw}",
            "testCode":   "import pytest\ndef test_fail(): pytest.fail('parse error')",
            "rationale":  f"Model output could not be parsed: {e}",
        }

# ── JSON-RPC models ────────────────────────────────────────────────────────
class MessagePart(BaseModel):
    kind: str; text: str

class Message(BaseModel):
    kind: str; messageId: str; role: str; parts: list[MessagePart]

class JsonRpcParams(BaseModel):
    message: Message

class JsonRpcRequest(BaseModel):
    jsonrpc: str = Field(..., pattern="^2\\.0$")
    method: str; params: JsonRpcParams; id: str

# ── Main endpoint ──────────────────────────────────────────────────────────
@app.post("/")
async def handle_message(request: Request) -> JSONResponse:
    body = await request.json()
    logger.info(f"Request id={body.get('id', 'unknown')}")

    try:
        rpc = JsonRpcRequest(**body)
    except Exception as e:
        return JSONResponse(status_code=400, content={
            "jsonrpc": "2.0", "id": body.get("id"),
            "error": {"code": -32600, "message": str(e)},
        })

    if rpc.method != "message/send":
        return JSONResponse(content={
            "jsonrpc": "2.0", "id": rpc.id,
            "error": {"code": -32601, "message": f"Unknown method: {rpc.method}"},
        })

    task_prompt = " ".join(
        p.text for p in rpc.params.message.parts if p.kind == "text"
    )

    t0 = time.perf_counter()
    solution = generate_solution(task_prompt)
    logger.info(f"Generated in {time.perf_counter()-t0:.2f}s")

    return JSONResponse(content={
        "jsonrpc": "2.0",
        "id": rpc.id,
        "result": {
            "kind": "message",
            "parts": [{"kind": "text", "text": json.dumps(solution, indent=2)}],
        },
    })

@app.get("/health")
async def health():
    return {"status": "ok", "model": MODEL_NAME, "alpha": STEERING_ALPHA}

if __name__ == "__main__":
    import uvicorn
    uvicorn.run("server:app", host="0.0.0.0", port=8080, workers=1)
\end{lstlisting}
\end{codebox}

\subsection{Deployment Notes for Lambda Labs H100}

\begin{codebox}
\begin{lstlisting}[style=bash]
# Install dependencies
pip install fastapi uvicorn transformers accelerate bitsandbytes torch

# Launch
python -m uvicorn server:app --host 0.0.0.0 --port 8080 --workers 1

# Single H100 (80 GB): use 4-bit quantisation
# BitsAndBytesConfig(load_in_4bit=True, bnb_4bit_compute_dtype=torch.bfloat16)
# Dual H100 (160 GB): BF16 70B fits natively with device_map="auto"
\end{lstlisting}
\end{codebox}

\subsection{Awareness of LogoMesh Scoring Internals}

The Green Agent scores Purple Agent output using:

\[
  \text{CIS} = (0.25R + 0.25A + 0.25T + 0.25L) \times \text{Red\_Penalty}
               \times \text{Intent\_Multiplier}
\]

Key implications for the steered model:

\begin{description}[leftmargin=2em, style=nextline]
  \item[\textbf{R (Rationale Integrity).}]
    Scored via cosine similarity using \ic{all-MiniLM-L6-v2}. The rationale
    must be semantically aligned with the task; generic boilerplate will score
    near zero.

  \item[\textbf{L (Logic Score).}]
    \textbf{Has no floor protection} (Architecture Overview \S5.1). This is
    the highest-variance component and the primary target for concept steering.
    The LLM scorer can hallucinate this score arbitrarily low. \emph{Note: Phase 2 Remediation Plans indicate this loophole will soon be closed. RepE concept vectors extracted prior to this patch will demonstrate maximal efficacy in demonstrating the flaw before the system is secured.}

  \item[\textbf{T (Testing Integrity).}]
    Formula: $T = 0.20 + 0.65 \times \text{test\_pass\_rate} +
    \text{specificity\_bonus}$. Tests run in a Docker sandbox with 128\,MB
    RAM, 50\% CPU, 50 PIDs. Keep tests lightweight --- avoid heavy imports.

  \item[\textbf{Red Agent.}]
    The MCTS engine (now restricted to \textbf{5 steps} by default, or 3 in
    \ic{fast\_lenient} mode, 2 branches, 60\,s limit) mutates \ic{sourceCode}
    to probe for vulnerabilities (now explicitly targeting \textbf{C/C++ and
    Solidity} for the Phase 2 Cybersecurity track).
    \emph{Note on Infrastructure:} The Red Agent's dynamic tool execution is
    being migrated to a \textbf{Persistent Sandbox (Sidecar)} via JSON-RPC to
    mitigate host compromise risks (the `Uroboros' threat). The Purple Agent's
    FastAPI JSON-RPC architecture described above provides the exact schematic
    needed for this IPC channel, demonstrating a dual-use architectural synergy
    between RepE steering infrastructure and Red Agent sandboxing. Steering toward correct, minimal code implicitly reduces the attack surface and limits the Red Agent's mutation opportunities.
\end{description}

\begin{bugbox}[Known Bug: DBOM Hash Mismatch]
  The Green Agent hashes the \emph{unsorted} JSON string (\ic{json.dumps(result)}),
  but stores the \emph{sorted} string (\ic{json.dumps(result, sort\_keys=True)}).
  Cryptographic verification of any database record against its DBOM hash will
  consistently fail. This is a pre-existing Green Agent bug, not a Purple Agent
  concern, but worth noting if you attempt audit-trail validation.
\end{bugbox}

% ─────────────────────────────────────────────────────────────────────────────
\clearpage
\section*{References}
\addcontentsline{toc}{section}{References}
% ─────────────────────────────────────────────────────────────────────────────

\begin{enumerate}[leftmargin=2em]
  \item Zou, A., Phan, L., Chen, S., et al.\ (2023).
        \textit{Representation Engineering: A Top-Down Approach to AI Transparency}.
        arXiv:2310.01405.

  \item Jha, S., et al.\ (2025).
        \textit{vec2vec: Encoder-Free Cross-Modal Representation Transfer via Supervised Vector Translation}.
        arXiv:2505.12540. NeurIPS 2025.

  \item LogoMesh Internal Documentation (2026).
        \textit{Architecture Overview: Phase 1 Onboarding}.
        Confidential.
\end{enumerate}

\vfill
\begin{center}
  \small\color{darkgray}
  \rule{8cm}{0.4pt}\\[4pt]
  Document generated: 2026-03-04 \quad|\quad LogoMesh Research \quad|\quad Confidential
\end{center}

\end{document}
